\documentclass[11pt,a4paper]{article}

\usepackage[utf8]{inputenc}
\usepackage[T1]{fontenc}
\usepackage[brazilian]{babel}
\usepackage{lmodern}
\usepackage[margin=2.5cm]{geometry}
\usepackage{amsmath,amssymb}
\usepackage{xcolor}
\usepackage[most]{tcolorbox}
\usepackage{enumitem}
\usepackage{hyperref}
\usepackage{fancyhdr}
\usepackage{titlesec}
\usepackage{microtype}
\usepackage{array}
\usepackage{booktabs}
\usepackage{listings}

\definecolor{deitelblue}{RGB}{31,73,125}
\definecolor{boapratcolor}{RGB}{0,112,60}
\definecolor{errocolor}{RGB}{178,34,34}
\definecolor{engcolor}{RGB}{31,73,125}
\definecolor{desempcolor}{RGB}{198,89,17}
\definecolor{portabcolor}{RGB}{102,46,145}
\definecolor{codebg}{RGB}{248,248,248}
\definecolor{codekeyword}{RGB}{0,0,180}
\definecolor{codestring}{RGB}{163,21,21}
\definecolor{codecomment}{RGB}{0,128,0}

\lstdefinestyle{cppcode}{
  language=C++,
  basicstyle=\ttfamily\small,
  keywordstyle=\color{codekeyword}\bfseries,
  stringstyle=\color{codestring},
  commentstyle=\color{codecomment}\itshape,
  numbers=left, numberstyle=\tiny\color{gray}, numbersep=8pt,
  backgroundcolor=\color{codebg}, frame=single, framesep=4pt,
  rulecolor=\color{deitelblue!50}, breaklines=true,
  showstringspaces=false, tabsize=4,
  morekeywords={string, size_t, cin, cout, endl, inline, template, typename,
                bool, true, false, nullptr, override, final, public, private,
                protected, explicit, friend, virtual, this, static, operator,
                dynamic_cast, typeid},
  literate={ã}{{\~a}}1 {á}{{\'a}}1 {é}{{\'e}}1 {ê}{{\^e}}1 {í}{{\'i}}1
           {ó}{{\'o}}1 {ô}{{\^o}}1 {õ}{{\~o}}1 {ú}{{\'u}}1 {ç}{{\c{c}}}1
}

\hypersetup{colorlinks=true, linkcolor=deitelblue, urlcolor=deitelblue,
  pdftitle={C: Como Programar - Cap. 21 - Autorrevisao}}

\pagestyle{fancy}
\fancyhf{}
\fancyhead[L]{\small\textit{C: Como Programar} --- Cap.~21}
\fancyhead[R]{\small Exerc\'icios de Autorrevis\~ao}
\fancyfoot[C]{\small\thepage}
\renewcommand{\headrulewidth}{0.4pt}

\titleformat{\section}{\Large\bfseries\color{deitelblue}}{\thesection}{1em}{}
\titleformat{\subsection}{\large\bfseries\color{deitelblue}}{\thesubsection}{1em}{}

\newtcolorbox{boaprat}{enhanced, breakable, colback=boapratcolor!8, colframe=boapratcolor,
  fonttitle=\bfseries, title={\textbf{Boa Pr\'atica de Programa\c{c}\~ao}},
  left=8pt, right=8pt, top=4pt, bottom=4pt}
\newtcolorbox{errocomum}{enhanced, breakable, colback=errocolor!8, colframe=errocolor,
  fonttitle=\bfseries, title={\textbf{Erro Comum de Programa\c{c}\~ao}},
  left=8pt, right=8pt, top=4pt, bottom=4pt}
\newtcolorbox{obseng}{enhanced, breakable, colback=engcolor!8, colframe=engcolor,
  fonttitle=\bfseries, title={\textbf{Observa\c{c}\~ao de Engenharia de Software}},
  left=8pt, right=8pt, top=4pt, bottom=4pt}
\newtcolorbox{resposta}{enhanced, breakable, colback=deitelblue!5, colframe=deitelblue!70,
  boxrule=0.6pt, left=8pt, right=8pt, top=4pt, bottom=4pt}

\title{
  \vspace{-1cm}
  {\Huge\bfseries\color{deitelblue} Cap\'itulo 21}\\[0.3em]
  {\Large\bfseries Programa\c{c}\~ao Orientada a Objetos: Polimorfismo}\\[0.8em]
  {\large\itshape Exerc\'icios de Autorrevis\~ao --- Resolu\c{c}\~ao Did\'atica}
}
\author{}
\date{}

\begin{document}
\maketitle
\thispagestyle{empty}
\vspace{-2em}

\begin{center}
\rule{0.85\textwidth}{0.5pt}\\[0.4em]
\textit{``Polimorfismo \'e o coroamento da OO em C++. Significa, literalmente, ``muitas formas'': uma s\'o chamada \texttt{p->earnings()} executa c\'odigo \emph{diferente} dependendo do tipo concreto do objeto apontado por \texttt{p}. Este cap\'itulo introduz tr\^es novidades inseparavelmente ligadas: fun\c{c}\~oes \texttt{virtual}, classes abstratas com fun\c{c}\~oes \texttt{virtual = 0}, e o despacho din\^amico via \emph{vtable}. Juntos, eliminam o c\'odigo \texttt{switch} repetitivo que polui programas n\~ao-OO.''}\\[0.4em]
\rule{0.85\textwidth}{0.5pt}
\end{center}

\vspace{1em}

% ============================================================
\section{Exerc\'icio 21.1 --- Vocabul\'ario do Cap\'itulo 21}
% ============================================================

\textit{Preencha os espa\c{c}os em cada uma das senten\c{c}as (a)--(i).}

\subsection*{(a) Tratar um objeto da classe-base como \_\_\_\_ pode causar erros}

\begin{resposta}\textbf{Resposta:} objeto da \textbf{classe derivada}.\end{resposta}

A dire\c{c}\~ao oposta da \emph{conversao implicita upcast} (derivada $\to$ base) \'e o \emph{downcast} (base $\to$ derivada), e ele \'e \emph{perigoso}: pode-se estar lidando com um objeto que n\~ao \'e realmente do tipo derivado, gerando comportamento indefinido. \texttt{dynamic\_cast} (Cap.~21.8) faz essa convers\~ao com seguran\c{c}a, retornando \texttt{nullptr} ou lan\c{c}ando exce\c{c}\~ao se o objeto n\~ao for compat\'ivel.

\subsection*{(b) Polimorfismo ajuda a eliminar a l\'ogica de \_\_\_\_}

\begin{resposta}\textbf{Resposta:} \texttt{switch}.\end{resposta}

C\'odigo procedural cl\'assico discriminaria tipos com \texttt{switch}: ``se for funcion\'ario assalariado, fa\c{c}a A; se for horista, fa\c{c}a B''. Polimorfismo substitui esse \texttt{switch} por uma chamada \emph{virtual} \'unica.

\begin{lstlisting}[style=cppcode,numbers=none]
// ESTILO PROCEDURAL (com switch):
switch (f.tipo) {
    case ASSALARIADO: pag = f.salario * 4; break;
    case HORISTA:     pag = calcularHorista(f); break;
    // ...
}

// ESTILO POLIMORFICO (sem switch):
pag = f.earnings();   // virtual dispatch
\end{lstlisting}

\subsection*{(c) Se uma classe cont\'em pelo menos uma fun\c{c}\~ao virtual pura, ela \'e \_\_\_\_}

\begin{resposta}\textbf{Resposta:} classe \textbf{abstrata}.\end{resposta}

\texttt{virtual void f() = 0;} declara uma \emph{fun\c{c}\~ao virtual pura}. Uma classe com pelo menos uma dessas \'e abstrata: \emph{n\~ao pode ser instanciada diretamente}. S\'o serve como classe-base para outras classes que implementem todas as fun\c{c}\~oes virtuais puras.

\subsection*{(d) Classes das quais objetos podem ser instanciados s\~ao \_\_\_\_}

\begin{resposta}\textbf{Resposta:} classes \textbf{concretas}.\end{resposta}

Uma classe \'e concreta se todas as suas fun\c{c}\~oes virtuais puras (herdadas das bases) foram implementadas.

\subsection*{(e) Operador para downcast seguro}

\begin{resposta}\textbf{Resposta:} \texttt{dynamic\_cast}.\end{resposta}

\begin{lstlisting}[style=cppcode,numbers=none]
Conta *p = obterConta();
ContaPoupanca *cp = dynamic_cast<ContaPoupanca*>(p);
if (cp) {
    // sucesso: p apontava para ContaPoupanca; agora cp aponta para o mesmo
    cp->calculaJuros();
}
// se p NAO apontava para uma ContaPoupanca: cp == nullptr
\end{lstlisting}

\texttt{dynamic\_cast} usa a informa\c{c}\~ao RTTI (\emph{Run-Time Type Information}) para checar se o objeto real \emph{\'e} do tipo desejado. Para ponteiros, retorna \texttt{nullptr} em caso de falha; para refer\^encias, lan\c{c}a \texttt{std::bad\_cast}.

\subsection*{(f) \texttt{typeid} retorna refer\^encia a \_\_\_\_}

\begin{resposta}\textbf{Resposta:} objeto \texttt{type\_info}.\end{resposta}

\begin{lstlisting}[style=cppcode,numbers=none]
#include <typeinfo>
Funcionario *f = obterFuncionario();
const std::type_info &t = typeid(*f);
cout << "Tipo real: " << t.name() << endl;
\end{lstlisting}

\texttt{type\_info::name()} retorna uma string com o nome do tipo (formato definido pela implementa\c{c}\~ao --- \texttt{g++} usa nome \emph{mangled}, \texttt{c++filt} pode demangle-lo).

\subsection*{(g) \_\_\_\_ envolve usar ponteiro base ou refer\^encia para chamar virtuais}

\begin{resposta}\textbf{Resposta:} \textbf{Polimorfismo} (polimorfismo de subtipo / din\^amico).\end{resposta}

A defini\c{c}\~ao operacional: \texttt{Base *p = new Derivada(); p->virt();} chama a vers\~ao de \emph{Derivada}, n\~ao a de \emph{Base}.

\subsection*{(h) Fun\c{c}\~oes que podem ser sobrepostas usam a palavra-chave \_\_\_\_}

\begin{resposta}\textbf{Resposta:} \texttt{virtual}.\end{resposta}

A palavra-chave \texttt{virtual} na declara\c{c}\~ao da base ativa o despacho din\^amico. As derivadas s\~ao livres para sobrep\^or (com ou sem a anota\c{c}\~ao \texttt{override} do C++11, recomend\'avel).

\subsection*{(i) Convers\~ao de ponteiro base $\to$ derivada \'e chamada de \_\_\_\_}

\begin{resposta}\textbf{Resposta:} \textbf{downcasting}.\end{resposta}

\textbf{Up} (subir na hierarquia: derivada $\to$ base) \'e sempre seguro e implicito. \textbf{Down} (descer: base $\to$ derivada) \'e arriscado e exige \texttt{dynamic\_cast} ou alternativa equivalente.

\begin{errocomum}
\textbf{Um downcast com \texttt{static\_cast} \'e sint\'aticamente v\'alido mas semanticamente perigoso:}
\begin{lstlisting}[style=cppcode,numbers=none]
Conta *p = new ContaCorrente(1000, 2);
// Programador acredita que p eh ContaPoupanca:
ContaPoupanca *cp = static_cast<ContaPoupanca*>(p);  // COMPILA!
cp->calculaJuros();   // COMPORTAMENTO INDEFINIDO
\end{lstlisting}
Use \texttt{dynamic\_cast} para qualquer downcast em hierarquias polim\'orficas.
\end{errocomum}

% ============================================================
\section{Exerc\'icio 21.2 --- Verdadeiro ou Falso}
% ============================================================

\subsection*{(a) Todas as virtuais em classe-base abstrata devem ser puras}

\begin{resposta}\textbf{Resposta:} \textbf{Falso}.\end{resposta}

\textbf{Corre\c{c}\~ao:} uma classe-base abstrata pode incluir fun\c{c}\~oes \texttt{virtual} \emph{com implementa\c{c}\~ao}, al\'em das puras. O que torna a classe abstrata \'e ter \emph{pelo menos uma} virtual pura, mas n\~ao \emph{todas} precisam s\^e-lo.

\begin{lstlisting}[style=cppcode,numbers=none]
class Funcionario {
public:
    virtual double earnings() const = 0;   // PURA
    virtual void print() const;            // VIRTUAL COM IMPLEMENTACAO
};
\end{lstlisting}

Nossa implementa\c{c}\~ao do 21.12 segue exatamente esse padr\~ao: \texttt{print} \'e \texttt{virtual} mas tem corpo na base; \texttt{earnings} \'e \texttt{= 0}.

\subsection*{(b) Referir-se a objeto da derivada com \emph{handle} da base \'e perigoso}

\begin{resposta}\textbf{Resposta:} \textbf{Falso}.\end{resposta}

\textbf{Corre\c{c}\~ao:} \emph{o contr\'ario \'e perigoso}: referir-se a um objeto da \emph{base} com \emph{handle} da \emph{derivada}. Quando temos \texttt{Base *p = new Derivada()}, \'e seguro --- \'e o cen\'ario normal do polimorfismo. O perigoso seria \texttt{Derivada *p = static\_cast<Derivada*>(umaBaseDeVerdade)}.

\subsection*{(c) Classe vira abstrata sendo declarada \texttt{virtual}}

\begin{resposta}\textbf{Resposta:} \textbf{Falso}.\end{resposta}

\textbf{Corre\c{c}\~ao:} \emph{classes} n\~ao s\~ao declaradas \texttt{virtual}. A palavra-chave \texttt{virtual} se aplica a \emph{fun\c{c}\~oes-membro} e \texttt{base classes} (na heran\c{c}a virtual, Cap.~24). Uma classe se torna abstrata ao conter pelo menos uma \emph{fun\c{c}\~ao} virtual pura.

\begin{lstlisting}[style=cppcode,numbers=none]
virtual class Forma { };  // ERRO de sintaxe

class Forma {              // OK: classe abstrata por causa do = 0
public:
    virtual void desenhar() const = 0;
};
\end{lstlisting}

\subsection*{(d) Se a base declara uma virtual, a derivada deve implementar para virar concreta}

\begin{resposta}\textbf{Resposta:} \textbf{Verdadeiro} (com nuance).\end{resposta}

A nuance: a regra vale quando a virtual da base \'e \emph{pura}. Se for s\'o \texttt{virtual} (sem \texttt{= 0}), a derivada \emph{n\~ao precisa} sobrescrev\^e-la: herda a vers\~ao da base. Mas pelo enunciado, claramente a quest\~ao se refere a virtuais puras.

\subsection*{(e) Programa\c{c}\~ao polim\'orfica pode eliminar a necessidade de \texttt{switch}}

\begin{resposta}\textbf{Resposta:} \textbf{Verdadeiro}.\end{resposta}

Esse \'e o \emph{benef\'icio prim\'ario} do polimorfismo, conforme discutiremos em detalhes nos exerc\'icios 21.3 e 21.4: substituir \texttt{switch} por chamadas virtuais elimina a maior fonte de manuten\c{c}\~ao e bugs em sistemas com m\'ultiplos tipos de entidade.

\begin{obseng}
N\~ao \'e que \texttt{switch} seja sempre ruim: para casos pequenos, com tipos enumerados pr\'oximos uns dos outros, \texttt{switch} pode ser mais leg\'ivel que polimorfismo. Mas, quando o n\'umero de casos cresce e a l\'ogica espec\'ifica de cada tipo se torna substancial, polimorfismo ganha duas frentes: organiza\c{c}\~ao (cada caso vira a classe correspondente) e extensibilidade (adicionar um caso n\~ao requer tocar em c\'odigo cliente).
\end{obseng}

% ============================================================
\section*{S\'intese conceitual}
% ============================================================

A autorrevisão estabeleceu os pilares do polimorfismo em C++:

\begin{itemize}[leftmargin=*,itemsep=0pt]
  \item \textbf{Mecanismo:} fun\c{c}\~oes \texttt{virtual} permitem que uma chamada via ponteiro/refer\^encia da base execute c\'odigo da derivada.
  \item \textbf{Abstratos vs concretos:} classes com pelo menos uma virtual pura s\~ao abstratas (n\~ao instanci\'aveis). Classes que implementam todas as virtuais puras s\~ao concretas.
  \item \textbf{Convers\~oes:} upcast \'e impl\'icito e seguro; downcast requer \texttt{dynamic\_cast}.
  \item \textbf{RTTI:} \texttt{typeid} e \texttt{type\_info} fornecem informa\c{c}\~oes de tipo em tempo de execu\c{c}\~ao.
  \item \textbf{Benef\'icio prim\'ario:} eliminar \texttt{switch} sobre tipos --- t\'ema dos exerc\'icios 21.3 e 21.4.
\end{itemize}

\vspace{1em}
\begin{center}\rule{0.5\textwidth}{0.4pt}\end{center}

\end{document}
